% Preamble
\documentclass[11pt, twoside]{report}

% Packages

\usepackage[utf8]{inputenc}
\usepackage[T1]{fontenc}                            %Sets fontenc to max level needed in file
\usepackage{graphicx}                               %To load and display images
\usepackage{caption}                                %For subfigures and cap env
\usepackage{subcaption}
\usepackage[a4paper, width=150mm,                   %Size of Paper
            top=25mm, bottom=25mm,
            bindingoffset=6mm]{geometry}
\usepackage{fancyhdr}                               %Header and Footers
\pagestyle{fancy}                                   %To disable for specific page use \pagestyle{<type>}
                                                    %<type>: plain for pagenr
                                                        %    empty for empty pagestyle
                                                        %    fancy to enable the default syle again
\fancyhead{}                                        %Resets the header
\fancyhead[RO,LE]{\chaptername~\thechapter}         %Header Definition: Chaptername and number: Right on Odd and Left on even pages
\fancyfoot{}                                        %Resets the footer
\fancyfoot[CE, CO]{\thepage}                        %Centered Page Number on odd and even pages
%Define output dir, because of error with: https://tex.stackexchange.com/questions/112953/error-when-using-minted-package-and-output-directory-option
%Cache can be removed normally, but in certain situations it can result in errors. Setup minted see: https://alipourmousavi.com/blog/index.php/2018/02/08/using-minted-package-in-latex-to-format-codes/
\usepackage[cache=false, outputdir=Z:/Users/jan22/CodeProjects/latex_template/auxil]{minted} %Must be before csquotes
\setminted{
    linenos=true,
    autogobble
}
\usepackage[english]{babel}                         %For splitting words
\usepackage{csquotes}                               %Used for babel/polyglossia and minted
\usepackage{setspace}                               %Set line spaceing
\setstretch{1.5}
\usepackage{amsmath}                                %Math package
\usepackage{color}                                  %Colored Links
\usepackage{hyperref}                               %For setting links
\hypersetup{
    colorlinks=false,                                %true, if links should be colorful
    linktoc=all,                                    %all => sections and subsections linked
    linkcolor=black,                                %Color of link
}
\usepackage{listings}                               %Environment for listings
%\usepackage{eurosym}                               %Euro sign
%Maybe change back to biber
\usepackage[style=authoryear,citestyle=authoryear,backend=bibtex]{biblatex}   % Bib
\addbibresource{main.bib}                           %Set bib file

%FOR minted
\renewcommand{\MintedPygmentize}{Z:/Users/jan22/AppData/Local/Programs/Python/Python37-32/Scripts/pygmentize} %Set path to script location
\renewcommand{\theFancyVerbLine}{\sffamily\textcolor[rgb]{0.5,0.5,0.5}{\scriptsize\arabic{FancyVerbLine}}}
% Create a new environment for breaking code listings across pages.
%Change lstlisting to figure if code fragments should be listed under figures
\newenvironment{longlisting}{\captionsetup{type=lstlisting}}{}
\renewcommand{\lstlistingname}{Code Fragment}       %Sets name of Listing space lstlisting - Will occured before each caption inside its environment
\renewcommand{\lstlistlistingname}{List of \lstlistingname s} %Sets title for \listoflstlistings




%CUSTOM VARS:
\newcommand{\submissionDay}{day-of-submission}
\newcommand{\submissionMonth}{month-of-submission}
\newcommand{\submissionYear}{year-of-submission}
\newcommand{\maintitle}{thesis-title}
\newcommand{\subtile}{thesis-title-addition}
\newcommand{\name}{author-title-name}
\newcommand{\supervisor}{supervisor-title-name}

% Document
\begin{document}
    % Titelblatt:
% \newpage\mbox{}\newpage
\cleardoublepage   % force output to a right page
\thispagestyle{empty}
\begin{titlepage}
    \begin{flushright}
        \includegraphics[width=0.4\linewidth]{../src/other-pages/title/Logo-A3.jpg}
    \end{flushright}
    %! Suppress = LineBreak
    \begin{flushleft}
        \vspace{0.5cm}

        \section*{\maintitle}
        %\subsection*{\subtile}

        \vspace{0.5cm}

        Bachelor Thesis 2\\
        In Partial Fulfillment
        of the Requirements for the Degree of
        Bachelor of Science

        \vspace{0.5cm}

        \textbf{Bachelor of Science in Software and Information Engineering}

        \vspace{1cm}
        University of Applied Sciences FH Vorarlberg\newline
        Software and Information Engineering

        \vspace{0.5cm}

        Supervised by:\newline
        \phantom{x}\hspace{3ex}\supervisor\newline
        \vspace{0.5cm}

        Submitted by:\newline
        \phantom{x}\hspace{3ex}\name\newline
        \newline
        Dornbirn, \submissionMonth\space\submissionYear
    \end{flushleft}
\end{titlepage}

\newpage
\addtocontents{toc}{\protect\setcounter{tocdepth}{-1}}
\chapter*{Abstract}\label{ch:abstract}
Text

\chapter*{AbstractGER}\label{ch:abstractger}
Text

\addtocontents{toc}{\protect\setcounter{tocdepth}{2}}
\newpage
%Table of contents
\tableofcontents
\renewcommand{\thechapter}{\Roman{chapter}}
\newpage

\setcounter{chapter}{0}
\addcontentsline{toc}{chapter}{\protect\numberline{I}List Of Abbreviations}
\chapter*{List Of Abbreviations}
\clearpage
\begingroup %Prevent with a group the creation of new pages of \listof<name>
\let\clearpage\relax
\addcontentsline{toc}{chapter}{\protect\numberline{II}List Of Figures}
\listoffigures
\addcontentsline{toc}{chapter}{\protect\numberline{III}List Of Tables}
\listoftables
\addcontentsline{toc}{chapter}{\protect\numberline{IV}List Of Code Fragments}
\lstlistoflistings
\endgroup
\cleardoublepage

%Content
\renewcommand{\thechapter}{\arabic{chapter}}
\setcounter{chapter}{0}
\chapter{Introduction Chapter}\label{ch:intro-chapter}
\section{Inversion of Control}\label{sec:inv-of-control}
The Inversion of Control pattern can be seen in many programming frameworks, often being a defining characteristic of them.
A programmer often controls when his code is called.
However, when using a framework, control is handed over to said framework.
This is Inversion of Control, also known as "The Hollywood Principle" (don't call us, we'll call you).
It is a key difference between a library and a framework, as in a library, you call the functions, while many frameworks invert the control.\parencite{fowler_bliki_2005}
\newline
\section{Dependency Injection}\label{sec:dep-injection}
A common issue in programming is how to wire together different elements.
If you have a team building a web controller and another team building a database backend, how do you get them to work together?
For Dependency Injection, the basic idea is that a separate object, an "assembler", populates a field in a class with the appropriate implementation of an interface.
There are three main ways of doing this: Constructor Injection, Setter Injection and Interface Injection.
\subsection{Constructor Injection}\label{subsec:cons-injection}
Like the name implies, the injection takes place in the constructor in such a case.
\begin{longlisting}
    \begin{minted}{java}
        public class MovieLister {
            private MovieFinder finder;
            public MovieLister(MovieFinder finder) {
                this.finder = finder;
            }
        }
    \end{minted}
    \caption[Dependency Injection Example]{A code example showing dependency injection via constructor.}
    \label{lst:ci-example}
\end{longlisting}
Note that MovieFinder is an interface in this example.
All classes that implement said interface can be passed into the constructor and are injected into the class that way.\parencite{fowler_inversion_2004}




\cleardoublepage
\renewcommand{\thechapter}{\Roman{chapter}}
\setcounter{chapter}{3}
%\pagenumbering{Roman}
\appendix
\addcontentsline{toc}{chapter}{\protect\numberline{V}Appendix}
\chapter*{Appendix}
%Bibliography
\setcounter{biburllcpenalty}{7000} %Adds breakpoints in URLs - https://tex.stackexchange.com/questions/134191/line-breaks-of-long-urls-in-biblatex-bibliography
\setcounter{biburlucpenalty}{8000}
\addcontentsline{toc}{chapter}{\protect\numberline{VI}Bibliography}
\printbibliography

%Eigenständigkeitserklärung
%\input{documents/content/resources/other-pages/6_eidesstatlicheerklaerung.tex}

\end{document}
